(1) 几何光学射线方程可从Fermat原理得 \(d[n\,\mathbf r']/ds=\nabla n\)。采用 \(z\) 作为近轴传播坐标且 \(|dx/dz|\ll1\)，可写 \(d(n x')/dz\simeq\partial n/\partial x\)。由 \(n^2(x)\simeq n_0^2(1-gx^2)\) 得 \(n\simeq n_0(1-gx^2/2)\)、\(\partial_x n\simeq-n_0gx\)，故 \(x''+gx=0\)。
(2) 这一常系数方程的解为 \(x(z)=x_0\cos(\sqrt g\,z)+(\theta_0/\sqrt g)\sin(\sqrt g\,z)\)。初始条件可由 \(z=0\) 的位置以及导数直接验证；横向轨迹是沿光纤轴的正弦振荡，光线被折射率向中心增大的梯度周期性拉回。
(3) 横向正弦振荡空间周期为 \(\Lambda=2\pi/\sqrt g\)。一族准直光线具有共同初始斜率 \(\theta_0=0\)、不同初始横向位置 \(x_0\)，轨迹同时在 \(\cos(\sqrt g\,z)=0\) 处穿过轴，因此首次共同聚焦距离为 \(z_f=\pi/(2\sqrt g)\)。
(4) 外侧折射率较低使偏轴光线虽然路程较长却在不同位置具有较高局部相速度，理想抛物分布可减小几何光线间的传播时差。然而当 \(gr^2\) 不小或光线不再近轴时，光线方程含非线性项，焦点不再严格重合，实际还须考虑波导模式与材料色散。
